% =====================================================================
% Add to preamble (one-time)
% =====================================================================
% \definecolor{promptbg}{RGB}{249,250,251}
% \definecolor{promptaccent}{RGB}{37,99,235}
%
% \lstdefinestyle{prettyprompt}{
%     basicstyle=\ttfamily\footnotesize,
%     backgroundcolor=\color{promptbg},
%     showstringspaces=false,
%     breaklines=true,
%     breakatwhitespace=true,
%     columns=fullflexible,
%     keepspaces=true,
%     frame=l,
%     framerule=2.5pt,
%     rulecolor=\color{promptaccent},
%     framesep=10pt,
%     xleftmargin=15pt,
%     xrightmargin=4pt,
%     aboveskip=14pt,
%     belowskip=14pt,
%     abovecaptionskip=2pt,
%     belowcaptionskip=8pt,
%     captionpos=t,
%     extendedchars=true,
%     inputencoding=utf8,
%     literate={—}{{\textemdash}}1
%             {–}{{\textendash}}1
%             {’}{{'}}1,
% }

% =====================================================================
% Task generation prompt
% =====================================================================
\subsection{Prompts for Task Synthesis}
\paragraph{Task generation prompt.}
\Cref{lst:taskgen_prompt} shows the prompt used by the task generator. Given the documentation of a sampled subset of APIs, a target difficulty / action-type / task-focus configuration, and a list of negative examples (already-generated tasks) and under-utilized APIs, the model produces \texttt{\{num\_tasks\}} candidate tasks per call. The configurable clauses (\texttt{\{difficulty\_clause\}}, \texttt{\{action\_type\_clause\}}, \texttt{\{task\_focus\_clause\}}) are sampled from the variants enumerated in \Cref{tab:taskgen_difficulty,tab:taskgen_action,tab:taskgen_focus}.

\begin{figure*}[t]
\begin{lstlisting}[style=prettyprompt, captionpos=b, caption={Task generation prompt.}, label={lst:taskgen_prompt}]
You are an AI model with API understanding capabilities. You are provided with documentation of various APIs supported by commonly used mobile phone apps within <tools></tools> XML tags below:
<tools>
{api_docs_context}
</tools>

Your job is to generate realistic tasks that users would ask a powerful coding-based tool calling agent to solve by utilizing the provided APIs. The agent is allowed to solve the task in an iterative multi-step fashion. In each step, the agent can produce python code (including if/else and loops if needed) that invokes one or multiple APIs. This python code will be executed using an environment with API execution support and a backend user database, and the execution results will be provided to the agent before the next step.

Carefully go through all the provided API documentations which include the input parameter descriptions and the output response schemas for each API. Then, generate {num_tasks} distinct tasks each satisfying all the below requirements:

# Task Requirements:
- **Solvability:** The task should be completely solvable by an agent using the provided APIs without any additional information from the user. You can assume that the agent already knows the user's name, username, email, phone number, birthday, sex, address, payment cards, login access tokens for all the apps, and current date/time.
- **Specificity:** Use specific details (names, amounts, dates etc) only where needed to make the task unambiguous. Do not include anything the agent can look up using the APIs.
- **Clarity:** The task should not have contradicting or confusing statements.
- **Interdependent subtasks:** The task SHOULD NOT be a simple concatenation of completely independent subtasks. Solving the task should require the agent to solve many interdependent subtasks.
- **Referential complexity** When possible, introduce indirect references to entities (people, objects, locations, amounts, etc.) as long as these references can be resolved completely using the given APIs.
- **Realism:** Write the task exactly as a real user would say it to a personal assistant — natural, high-level, and goal-oriented. The task MUST NOT read like step-by-step instructions. The task should only state what the user wants and the agent will figure out how to achieve it.
- **Coverage:** Each task should cover {min_apis_per_app}-{max_apis_per_app} apis from each app.
- **Required apps:** Every task MUST use at least one of the following apps: {required_apps}. Tasks that do not involve any of these apps are invalid and will be rejected.The APIs covered by each task should vary between the tasks.
- **Difficulty:** {difficulty_clause}
- **Action type:** {action_type_clause}
- **Task focus:** {task_focus_clause}

{focus_section}
{negative_section}
First carefully think about how to construct tasks that satisfy all the above requirements. Specifically, think about which subsets of the provided APIs can result in meaningful tasks and identify several such subsets to generate multiple tasks. Do not make any assumptions about any API. The API documentation clearly describes the input parameters and the expected output. DO NOT hallucinate any new APIs that are not part of the APIs provided within <tools></tools> tags. For each task, first generate the names of all the APIs (comma separated) covered by the task in a separate line starting with `APIs COVERED:`, then describe how an agent can fully solve the task using these APIs in a separate line starting with `EXPLANATION:`, and finally generate the task in a separate line that starts with `TASK:`. You should not include any additional information in this line other than the task itself. DO NOT generate any additional summaries or titles for the tasks.
\end{lstlisting}
\end{figure*}


% =====================================================================
% Difficulty levels
% =====================================================================
\begin{table*}[t]
\centering
\small
\renewcommand{\arraystretch}{1.3}
\begin{tabular}{l p{0.85\linewidth}}
\toprule
\textbf{Level} & \textbf{Clause} \\
\midrule
\textsc{easy}   & The task requires the agent to perform a focused, self-contained operation. The solution would be a linear sequence of API calls with minimal branching. Filtering, sorting, or iterating over items would be the main challenge. \\
\midrule
\textsc{medium} & The task introduces indirect references that the agent must resolve via API calls before acting (e.g.\ ``the person I owe the most'', ``my most recent group'', ``the highest-rated item''). The output of one subtask feeds directly into the next, creating a clear dependency chain. The solution requires moderate branching or iteration. \\
\midrule
\textsc{hard}   & The task requires multi-hop reasoning: the agent must gather context across multiple apps, reconcile entities, and execute a coordinated sequence of actions. For example, an entity resolved in app A feeds a lookup in app B which feeds an action in app C. The solution requires conditional branches or loops over dynamic lists. \\
\bottomrule
\end{tabular}
\caption{Difficulty clauses inserted into the task generation prompt. Exactly one is sampled per task.}
\label{tab:taskgen_difficulty}
\end{table*}


% =====================================================================
% Action type
% =====================================================================
\begin{table*}[t]
\centering
\small
\renewcommand{\arraystretch}{1.3}
\begin{tabular}{l p{0.85\linewidth}}
\toprule
\textbf{Type} & \textbf{Clause} \\
\midrule
\textsc{mixed} & The task should be a MIXED task that requires substantial reading AND writing: the agent must first gather and reason over information from one or more apps (\texttt{GET} requests), then use that information to drive meaningful write operations (\texttt{POST}, \texttt{PATCH}, \texttt{DELETE}) in the same or different apps. Neither reads nor writes dominate --- both are essential to completing the task. \\
\midrule
\textsc{read}  & The task should be a READ-HEAVY task. The task primarily requires the agent to query, search, look up, and report information. The majority of API calls should be \texttt{GET} requests (lookup, list, search, show). Write operations should be minimal or absent --- the output of the task is a piece of information or a summary, not a state change in any app. \\
\midrule
\textsc{write} & The task should be a WRITE-HEAVY task. The task primarily requires the agent to create, update, send, buy, delete, or modify data. The majority of API calls should be \texttt{POST}, \texttt{PATCH}, or \texttt{DELETE} requests. Read operations are only used minimally to look up IDs or verify preconditions before writing. \\
\bottomrule
\end{tabular}
\caption{Action-type clauses inserted into the task generation prompt. Exactly one is sampled per task.}
\label{tab:taskgen_action}
\end{table*}


% =====================================================================
% Task focus
% =====================================================================
\begin{table*}[t]
\centering
\small
\renewcommand{\arraystretch}{1.3}
\resizebox{1\linewidth}{!}{%
\begin{tabular}{l p{0.82\linewidth}}
\toprule
\textbf{Focus} & \textbf{Clause} \\
\midrule
\textsc{constraint\ satisfaction} & The task should require the agent to find or select entities that satisfy MULTIPLE simultaneous criteria. The agent must filter, search, or verify that several conditions hold at the same time before acting. Examples of constraints: price within a range,  rating above a threshold, date before a deadline, belonging to a specific category, matching a text pattern, being associated with a particular person. The task should require at least 2--3 distinct constraints to be checked together. The challenge is in correctly combining the constraints, not in the action itself. \\
\midrule
\textsc{derivation} & The task should require the agent to derive a SINGLE piece of new information by aggregating or comparing multiple data points, then act on that derived result. The derived value is essential --- the agent cannot complete the task without first computing it. Examples: compute total spending to decide a budget, find who owes the most and message that one person, calculate average rating to set a threshold, determine which playlist has the longest total duration and play it.
The task must involve an aggregation step (sum, average, max, min, count, ranking, or comparison) whose SINGLE result drives a subsequent action. The final action should target one entity or a fixed number of entities --- NOT loop over every item in a set. \\
\midrule
\textsc{iteration} & The task should require the agent to perform an operation on EACH item in a set whose membership is discovered at runtime. The set is not known in advance --- the agent must first query to find which items qualify, then process each one individually. Each iteration may involve multiple API calls. Examples: for each person who owes me  money send a reminder, for every overdue task add a comment, for all notes tagged `urgent' share them with a colleague, for each expense above \$50 update its description. The task should NOT be solvable with a single bulk API call --- individual per-item processing is required. \\
\midrule
\textsc{open} & \emph{(no additional focus constraint)} \\
\bottomrule
\end{tabular}%
}
\caption{Task-focus clauses inserted into the task generation prompt. Exactly one is sampled per task. The \textsc{open} setting adds no additional constraint.}
\label{tab:taskgen_focus}
\end{table*}


% =====================================================================
% Task judge prompt
% =====================================================================
\paragraph{Task judge prompt.}
After a candidate task is generated, a judge LLM verifies that it satisfies the same requirements that were given to the generator. \Cref{lst:judge_prompt} shows the full prompt template.

\begin{figure*}[t]
\begin{lstlisting}[style=prettyprompt, captionpos=b, caption={Task judge prompt.}, label={lst:judge_prompt}]
Your job is to evaluate the quality of agentic tasks generated by a task generator. The task generator creates tasks based on the descriptions of various APIs provided to it.

Here are the API descriptions provided to the task generator:
{api_docs_context}

Here are the task requirements given to the task generator:
{task_requirements}
Evaluate if the task satisfies all the requirements. In order to evaluate, you should briefly think about how to solve the task. This will help you to figure out the API coverage, the task difficultly and its focus. DO NOT THINK FOR TOO LONG. First provide a brief reasoning behind your decision and then provide your final YES/NO answer in a separate line that starts with `ANSWER:`. You should not include any additional information in this line other than the YES/NO answer itself.

Here is the task to evaluate:
{task_description}
\end{lstlisting}
\end{figure*}


% =====================================================================
% Task rewriting prompt — body
% =====================================================================
\paragraph{Task rewriting prompt.}
Tasks that pass the judge are then rewritten to remove procedural scaffolding and sound like a natural user request. \Cref{lst:rewrite_prompt} shows the prompt body, and \Cref{lst:rewrite_examples} shows the three in-context examples appended after it.

\begin{figure*}[t]
\begin{lstlisting}[style=prettyprompt, captionpos=b, caption={Task rewriting prompt (body).}, label={lst:rewrite_prompt}]
You are an AI model with API understanding capabilities. Your job is to generate realistic tasks that users would ask a powerful coding-based tool calling agent to solve by utilizing a set of available APIs. The agent is allowed to solve the task in an iterative multi-step fashion. In each step, the agent can produce python code (including if/else and loops if needed) that invokes one or multiple APIs. This python code will be executed using an environment with API execution support and a backend user database, and the execution results will be provided to the agent before the next step.

You will be provided with a task and the documentations of APIs that are needed for solving this task. Carefully go through all the provided API documentations which include the input parameter descriptions and the output response schemas for each API. Then, rewrite the given task in a few sentences such that the rewritten task satisfies all the below requirements:

# Task Requirements:
{task_requirements}
Make sure that the rewritten task also requires all the APIs needed for solving the original task. Do not add any new information or change the intent/scope of the task. Make sure to remove redundant setup phrases, obvious intermediate steps that any agent would infer, and overly explicit instruction-style wording that describes the procedure rather than the goal. Make sure that you do not remove any details that are crucial to what the user exactly wants. Examples of things that must be kept:
    * Specific file names, folder paths, or formats (e.g. 'save it as a .txt file', 'store in ~/reports/', 'delete the zip file after')
    * Specific actions the user explicitly wants (e.g. 'then delete it', 'send them a message', 'create a note')
    * Specific names, amounts, dates, or identifiers
    * Any condition or constraint the user cares about (e.g. 'only if they haven't texted me')

First provide a brief explanation of your rewriting process. Then provide the list of all the APIs (comma separated) covered by the rewritten task in a separate line starting with `APIs COVERED:`, and finally provide the rewritten task in a separate line that starts with `TASK:`. You should not include any additional information in this line other than the task itself.

Here is the task to rewrite:
{task_description}

Here are the APIs required by this task:
{api_docs_context}
\end{lstlisting}
\end{figure*}


% =====================================================================
% Task rewriting prompt — in-context examples
% =====================================================================
\begin{figure*}[t]
\begin{lstlisting}[style=prettyprompt, captionpos=b, caption={Task rewriting prompt (in-context examples appended to the body).}, label={lst:rewrite_examples}]
Here are some in-context examples to guide you with the rewriting process:
## Example 1:
Original: I want to organize sharing photos with my family. Find all my family contacts and check whether I have created a dedicated folder for each person named `<first_name>_<last_name>` in `~/photos/`. If they do have a folder, determine whether we have ever exchanged text messages. For any family member who has a photo folder but whom I have never texted, create a note titled `<first_name> <last_name> Photos` to remind me that I need to share photos with them.

Rewritten: For any of my family contacts who have a folder named `<first_name>_<last_name>` in `~/photos/` but whom I have never texted, create a note titled `<first_name> <last_name> Photos` stating that I need to share photos with them.

## Example 2:
Original: I need to check up on the people listed in my `Project Alpha Team` note. They are supposed to submit their reports as text files named `<first_name>_<last_name>.txt` in the `~/project_alpha/reports/` folder. Determine who hasn't submitted their report yet. For those whose reports are missing, check whether they have sent me any text messages recently. Provide a list of the team members who are both missing their reports and haven't been communicating via text.

Rewritten: Who are the team members listed in my `Project Alpha Team` note who haven't submitted their reports named `<first_name>_<last_name>.txt` to the `~/project_alpha/reports/` folder and also haven't texted me recently?

## Example 3:
Original: Download the latest invoice from my email, save it to ~/invoices/latest.pdf, and then delete the email.

Rewritten: Save the latest invoice from my email to ~/invoices/latest.pdf, then delete that email.
\end{lstlisting}
\end{figure*}
